%% LyX 2.4.4 created this file.  For more info, see https://www.lyx.org/.
%% Do not edit unless you really know what you are doing.
\documentclass[11pt,oneside,english]{amsart}
\usepackage{mathptmx}
\usepackage[LGR,T1]{fontenc}
\usepackage{textcomp}
\usepackage[latin9]{inputenc}
\usepackage{babel}
\usepackage{amsthm}
\usepackage{geometry}
\geometry{verbose}
\usepackage{jurabib}[2004/01/25]

\makeatletter

%%%%%%%%%%%%%%%%%%%%%%%%%%%%%% LyX specific LaTeX commands.

\newcommand*\LyXZeroWidthSpace{\hspace{0pt}}
\DeclareRobustCommand{\greektext}{%
  \fontencoding{LGR}\selectfont\def\encodingdefault{LGR}}
\DeclareRobustCommand{\textgreek}[1]{\leavevmode{\greektext #1}}


%%%%%%%%%%%%%%%%%%%%%%%%%%%%%% Textclass specific LaTeX commands.
\numberwithin{equation}{section}
\numberwithin{figure}{section}

\AtBeginDocument{
  \def\labelitemi{\(\ast\)}
}

\makeatother

\begin{document}

\section{Overview }

\subsection{Objective}

We estimate the capitalization effect of local sewage spill exposure
on house prices.

The key empirical challenge is that CSO spill behaviour is highly
persistent. Some CSOs spill frequently and others rarely do, meaning
most variation is cross-sectional. Houses near chronically bad CSOs
are always cheaper. This makes it difficult to distinguish spill effects
from unobserved neighbourhood quality.

Our strategy generates plausibly exogenous within-area, over-time
variation in spill exposure using rainfall shocks and predetermined
hydraulic constraints.

\subsection{Identification Logic}

The strategy combines three components:
\begin{enumerate}
\item Rainfall shocks (exogenous time variation)
\begin{enumerate}
\item Rainfall provides plausibly exogenous shocks to sewer systems. 
\item On its own, rainfall affects all areas, but it does not translate
uniformly into spills.
\end{enumerate}
\item Predetermined hydraulic constraints (supply-side heterogeneity). Whether
rainfall causes spills depends on hydraulic stress:
\begin{enumerate}
\item Wastewater treatment works (WwTW) loading relative to design capacity
\item CSO baseline hydraulic sensitivity. Or, elevation-based hydraulic
slack within the interceptor system
\item These constraints are predetermined by historical engineering design
and long-run infrastructure capacity.
\item Rainfall \texttimes{} hydraulic stress generates differential spill
responses across systems.
\end{enumerate}
\item River-network exposure (demand-side mapping). 
\begin{enumerate}
\item Once a CSO spills, sewage enters the river network and travels downstream. 
\item Exposure depends on:
\begin{enumerate}
\item Whether a house lies downstream of the CSO,
\item Distance along the river network,
\item Tidal conditions.
\end{enumerate}
\item We construct house-level spill exposure by aggregating CSO-level spills
using downstream river-distance weights.
\end{enumerate}
\end{enumerate}

\subsection{Empirical Implementation}

\textbf{Step 1: First Stage (CSO-level)}

We estimate how rainfall shocks interacting with predetermined hydraulic
stress affect CSO spill intensity, including CSO fixed effects and
time fixed effects.

Identification comes from differential translation of rainfall into
spills across systems with different hydraulic constraints --- not
from baseline spill levels.\\

\noindent\textbf{Step 2: Exposure Mapping}

Predicted CSO spills are aggregated to the house level using river-network
connectivity weights. Only downstream exposure enters the baseline
specification.\\

\noindent\textbf{Step 3: House Price Equation}

We estimate:

{*} Log house price on recent downstream spill exposure,

{*} Controlling flexibly for rainfall in the same window,

{*} Including neighbourhood fixed effects and quarter-time fixed effects.

Identification comes from within-neighbourhood, within-time variation
in spill exposure induced by rainfall \texttimes{} hydraulic constraints.

Rainfall is controlled for directly, so the instrument operates only
through the spill channel.

\subsection{Exclusion Restriction}

Conditional on neighbourhood fixed effects and rainfall controls:

{*} Hydraulic stress affects house prices only through CSO spill intensity.

{*} It does not directly alter flood risk, schools, transport, or
other amenities.

{*} River-network exposure is mechanically determined by geography. 

\subsection{Alternative Instruments and Robustness}

We consider:

\textbf{Alternative instruments}

{*} Hydraulic slack (elevation gradient relative to WwTW),

{*} Soil permeability (geological runoff capacity),

{*} WwTW loading ratios.

These are predetermined engineering or geological features and enter
only through interaction with rainfall shocks.

\textbf{Time Aggregation}

Buyers may capitalize long-run spill expectations rather than single
events.

We therefore examine:

{*} Short-run spill exposure windows (salience effects),

{*} Annual or rolling averages (long-run expectations).

Short-run windows provide sharper quasi-experimental variation; longer
windows capture belief formation.

\section{Baseline }

We exploit exogenous rainfall shocks interacting with predetermined
hydraulic constraints at the wastewater treatment works and CSO level
to generate differential variation in spill intensity. House-level
exposure varies mechanically according to downstream river connectivity.
Within-neighbourhood, within-time variation in predicted spill exposure
identifies the capitalization effect.

Our identification combines (i) exogenous rainfall shocks, (ii) predetermined
hydraulic constraints that govern how rainfall translates into CSO
spills, and (iii) river-network connectivity that determines which
houses are exposed once spills occur. The strategy proceeds in three
steps: first, we estimate how rainfall interacting with hydraulic
stress predicts CSO-level spill intensity; second, we map predicted
spills to houses using downstream river connectivity; third, we estimate
the capitalization of this predicted spill exposure in house prices
within neighbourhood and time fixed effects.

Rainfall shocks \texttimes{} predetermined hydraulic constraints \textrightarrow{}
differential spill intensity \textrightarrow{} river-based house exposure
\textrightarrow{} capitalization in house prices.

Notation
\begin{itemize}
\item house $h$, quarter $q$, year $t$ in MSOA $m\left(h\right)$
\item CSO $c$ in Waste-water Treatment Works catchment area $w\left(c\right)$ 
\item $\alpha$ are fixed effects
\begin{itemize}
\item $m\left(h\right)$ is neighbourhood $m$ of house $h$ (that could
be defined as catchment area, MSOA or the house itself at the very
finest)
\item $m\left(c\right)$ is the local area $m$ if CSO $c$ (that is defined
as the wastewater treatment works catchment area because it keeps
the sewage system across CSOs constant)
\end{itemize}
\item $\varepsilon$ and $\eta$ are usual random error terms 
\end{itemize}
We are interested in identifying the capitalization effect of local
sewage spill exposure on house prices. The key econometric challenge
is that spills are highly persistent at a given CSO --- meaning some
CSOs consistently spill a lot and others rarely do. Most of the identifying
variation is cross-sectional and we would expect: houses near bad
CSOs are always cheaper, houses near good CSOs are always more expensive.
With transaction-level data and house fixed effects, all that cross-sectional
variation is wiped out, and the effect of sewage spills on house prices
\textgreek{\textbeta} is identified only from within-CSO changes in
spill intensity over time, which may be limited. But without house
fixed effects, there is a risk of confounding neighbourhood quality
with sewage spills --- the CSOs that spill most are probably in older,
denser, poorer urban areas. 

To address this, the empirical strategy exploits exogenous rainfall
shocks interacting with predetermined wastewater treatment works capacity
constraints to generate differential variation in CSO spill intensity
that then translate into heterogeneous house-level exposure depending
on downstream river connectivity of the house to each CSO. The main
instrument is the interaction with stress on the wastewater treatment
works and a second instrument is the interaction with baseline CSO
spill propensity, that will also be discussed later in the alternative
section. 

The Waste-water Treatment Works (WwTW) catchment boundary is relevant
for the supply side of spills --- whether a CSO activates at all.
When a combined sewer surcharges and the WwTW is at capacity, dry
weather flow can no longer be routed to treatment, and the CSO opens
to relieve pressure. So the propensity to spill is a function of the
whole catchment's hydraulic load relative to the WwTW's treatment
capacity. We measure this with stress on the WwTW, along with other
stress factors such as catchment imperviousness and sewer vintage.
CSO spilling is specified as 
\begin{alignat*}{1}
Spill_{ct}=\pi_{1}\LyXZeroWidthSpace(R_{ct}\LyXZeroWidthSpace\times Stress_{w0}\LyXZeroWidthSpace)+\pi_{2}\LyXZeroWidthSpace(R_{ct}\LyXZeroWidthSpace\times Spillc_{c0})+\pi_{3}\LyXZeroWidthSpace R_{ct}\LyXZeroWidthSpace+\alpha_{c}\LyXZeroWidthSpace+\alpha_{t}\LyXZeroWidthSpace+\eta_{ct}
\end{alignat*}
where $0$ denotes initial period mean and rainfall $R_{ct}$ is the
time-varying shock. The inflow arriving at a specific CSO node depends
primarily on rainfall in its upstream drainage area. This is spatially
local to the CSO. Even if local inflow rises, whether the CSO spills
depends on how full the interceptor sewer is, whether upstream CSOs
already relieved pressure. If the WwTW is near or above dry-weather
capacity, additional inflow cannot be treated, interceptor surcharge
rises, CSOs open to relieve pressure. This part is catchment-wide.
The aggregate hydraulic burden on the treatment work is
\[
Stress_{w0}=ActualPE_{w0}/DesignPE_{w0},
\]
the load on the wastewater treatment works relative to the designated
load that it can carry (see more in the WwTW sub-section later). The
interaction term captures differential translation of rainfall shocks
into spill intensity across systems with different predetermined capacity
constraints, not just precipitation. It accounts for spatial variation
in runoff coefficients that rainfall alone does not, and is directly
measured at the works, and not just inferred from a weather station. 

$Spill_{c0}$ is the initial spill frequency of the CSO that captures
its local hydraulic signature. Because baseline propensity and stress
may be correlated with neighbourhood characteristics, CSO fixed effects
are included, so that only the interaction matters. Identification
comes from how rainfall shocks translate into spill intensity differentially
across systems with different predetermined capacity stress, not from
baseline levels.

The watercourse or river network is relevant for the demand side of
exposure --- which houses are actually affected once a CSO spills.
The discharged sewage enters the watercourse and flows downstream,
so exposure is a function of where the house is relative to the receiving
water body, not where it is relative to the WwTW. Houses along the
river downstream of a CSO outfall are exposed; houses in the same
WwTW catchment but draining to a different watercourse may not be
at all.

We construct a house-specific spill exposure measure by aggregating
CSO-level spills using river-network distance weights. How much a
spill at a given CSO affects a house is primarily through its distance
along the river network from the CSO outfall point to the closest
point on river near the house, $d_{hc}$. This can differ based on
downstream-upstream direction in the directed river graph, denoted
by an indicator $Downstream_{hc}$. Upstream impacts from a point
discharge are usually limited to extreme storm events and to rivers
with tidal reversals that make backflows feasible. We account for
this and define exposure $f_{hc}$ flexibly as a function of these
characteristics to obtain the house-specific spill variable as a sum
over all relevant CSOs that the house could be exposed to:

\begin{alignat*}{1}
Spill_{ht}=\sum_{c\in C\left(h\right)}\underbrace{f\left(d_{hc},Downstream_{hc},Tidal_{hc}\right)}_{f_{hc}}Spill_{ct}
\end{alignat*}
The function $f_{hc}$ can be defined as different distance bins that
are allowed to vary by downstream-upstream and tidal-non tidal rivers
or parametrically as $Downstream_{hc}\times\exp\left(-d_{hc}/\kappa\right)\times\left(1-Tidal_{hc}\right)$
for an exponential decay function (with decay parameter $\kappa$)
and $Downstream_{hc}\times\left(1+d_{hc}/\kappa\right)^{-\rho}\times\left(1-Tidal_{hc}\right)$
for the usual power law form. Similarly, $C\left(h\right)$ can refer
to all CSOs in the MSOA or to a subset that are within a certain river
distance or a certain upstream river distance of the house. 

Within the same micro-area and same time, the house price equation
is specified as follows to capture whether houses sold after more
recent nearby downstream CSO discharge sell for less/more, holding
rainfall conditions fixed:
\begin{alignat*}{1}
\ln Price_{ht}= & \beta\sum_{k=1}^{L}Spill_{h,t-k}+\sum_{k=1}^{L}\pi_{k}R_{h,t-k}+\Gamma\prime X_{ht}\LyXZeroWidthSpace+\alpha_{m(h)}\LyXZeroWidthSpace+\alpha_{qt}\LyXZeroWidthSpace+\varepsilon_{ht}
\end{alignat*}
The predicted CSO-level spills from the first stage are aggregated
to the house level using $f_{hc}$ weights to form the second stage
regressor $Spill_{ht}$. Identification comes from within-neighbourhood,
within-quarter variation in recent spill exposure. The coefficient
\textgreek{\textbeta} captures the semi-elasticity of house prices
with respect to recent downstream spill exposure.

Rainfall $R_{h,t-k}$ in the $L$-period window is included to allow
for direct effects. Rainfall affects house prices through flood risk
and amenity directly, and partialling it out means the instrument
only works through the spill channel. Because rainfall in the same
window is included, identification comes from variation in spill exposure
induced by rainfall \texttimes{} hydraulic constraints, beyond any
direct effect of rainfall on house prices.

$X_{ht}$ consists of house attributes (type, size, bedrooms, new
build, tenure, EPC band, etc.) and other area characteristics. Along
with the controls, we include neighbourhood (MSOA/OA) fixed effects
$\alpha_{m\left(h\right)}$ and $\alpha_{qt}$ are quarter-time fixed
effects to compare house prices within neighbourhoods and in the same
time period. 

Alternatively, the two steps can be combined by defining $Spill_{ht}=\sum_{c\in C\left(h\right)}\underbrace{f\left(d_{hc},Downstream_{hc},Tidal_{hc}\right)}_{f_{hc}}Spill_{ct}$
and clustering standard errors at the WwTW level to estimate the following
system of equations:

\begin{alignat*}{1}
\ln Price_{ht}= & \beta Spill_{ht}+\pi R_{ht}+\Gamma\prime X_{ht}\LyXZeroWidthSpace+\alpha_{m(h)}\LyXZeroWidthSpace+\alpha_{qt}\LyXZeroWidthSpace+\varepsilon_{ht}\\
Spill_{ht}= & \delta_{1}\sum_{c\in C\left(h\right)}f_{hc}\left(R_{ct}\times Stress_{w0}\right)+\delta_{2}\sum_{c\in C\left(h\right)}f_{hc}\left(R_{ct}\times Spill_{c0}\right)\\
 & +\zeta R_{ht}+\Upsilon\prime X_{ht}+\alpha_{m(h)}+\alpha_{qt}+\eta_{ht}
\end{alignat*}
Given that the instruments vary at the WwTW/CSO level but the regression
is at the house level, standard errors need to be cluster at the WwTW/CSO
level at minimum. With spatial data, considering whether there is
residual spatial autocorrelation across nearby WwTW catchments that
requires two-way clustering or Conley standard errors.

Conditional on neighbourhood fixed effects and rainfall controls,
the interaction between rainfall shocks and predetermined hydraulic
stress affects house prices only through CSO spill intensity. Hydraulic
stress does not directly alter flood risk, school quality, transport
access, or other amenities. Rainfall itself is controlled flexibly.
Once we have a strong first stage, the spatial correlation can also
be explored in more detail. Rainfall shocks are spatially correlated.
River networks connect across WwTWs. Housing markets spill over.

\section{Alternative}

\subsubsection*{CSO Hydraulics}

The baseline uses $Spill_{c0}$ because it is likely to be a strong
predictor of future spills. But baseline spill propensity reflects
long-run infrastructure quality, which correlates with neighbourhood
quality. To strengthen the instrument, we can also consider interactions
of rainfall shocks with CSO-level impervious surface, soil impermeability,
elevation relative to river and WwTW, distance to WwTW and density
of CSOs in the local area, that are more underlying determinants of
$Spill_{c0}$. 

\subsubsection*{Elevation-based hydraulic slack}

The main hydraulically-determined characteristic that determines how
much flow a CSO can pass before spilling is CSO elevation relative
to the treatment works that it conveys the flows into. Let $c$ denote
a CSO, belonging to waste treatment plant catchment area $w$. Ground
elevation at CSO location is $z_{c}$ and at the treatment works inlet
is $z_{w}$, with the distance between the CSO and the treatment inlet
given by $L_{cw}$. A CSO at high elevation relative to the treatment
works has a steep available gradient that can more readily achieve
self-cleansing velocity relative to a CSO at low elevation (near treatment
works elevation) that has a shallow available gradient. High elevation
CSOs have greater propensity to convey flows onwards to the treatment
works without spilling. Flow from a CSO to the treatment works travels
through the interceptor sewer. The interceptor must maintain a minimum
gradient to flow by gravity. The elevation difference between the
CSO and the treatment works, divided by the pipe distance between
them, determines the available hydraulic gradient $\left(z_{c}-z_{w}\right)/L_{cw}$. 

The hydraulic slack at each CSO --- how much spare capacity the interceptor
has at the point where the CSO's feeder meets it is $S_{cw}-S_{\min}\equiv\left(z_{c}-z_{w}\right)/L_{cw}-S_{\min}$where
$S_{\min}$ is the minimum slope needed to achieve self-cleansing
velocity (\ensuremath{\approx} 1/500 for a typical Victorian combined
sewer). Within a treatment works catchment, a CSO that sits 20 metres
higher than the works has a fundamentally different hydraulic situation
from one that sits at the same elevation. $S_{cw}$ is positive when
the available gradient exceeds the minimum needed to convey flow ---
i.e., the system has hydraulic slack. It is negative when the gradient
is insufficient --- the system is already constrained, and pumping
or surcharge is likely. The instrument normalises each CSO's elevation
disadvantage relative to the worst-positioned CSO in the catchment
of a treatment works: 
\begin{alignat*}{1}
E_{cw}\equiv & \max\left(0,S_{\min}-S_{cw}\right)/\max_{c\in w\left(c\right)}\left(\max\left(0,S_{\min}-S_{cw}\right)\right)
\end{alignat*}

Within-treatment-works fixed effects absorb the historical siting
of the works; identification comes from cross-CSO variation in relative
elevation within the same sewer system. A CSO with positive slack
$S_{cw}>S_{\min}$ has $E_{cw}=0$; a CSO with negative slack has
$E_{cw}>0$, proportional to how hydraulically constrained it is.
The variation in the instrument is hydraulically meaningful because
it determines whether additional flow causes a CSO to tip over into
spilling. It is pre-determined by Victorian interceptor routing (terrain-constrained)
and treatment works location (historically fixed). It can be measured
from OS Terrain 5 DEM ($z_{c},z_{w}$) and Hoffmann et al. (2025)
catchment data ($L_{cw}$ approximated by straight-line \texttimes{}
1.3 urban tortuosity factor, with sensitivity for 1.2-1.4). 

Conditional on treatment works catchment fixed effects, the elevation
of a specific CSO location relative to the treatment works elevation
is determined by the topography of the terrain (geological), the historical
decision about where to site the treatment works (pre-modern), and
where the Victorian interceptor was routed (pre-modern, terrain-constrained).
CSO elevation relative to the works may correlate with floodplain
or valley location; this is absorbed by micro-area fixed effects in
the price equation and WwTW fixed effects in the first stage. Additionally,
we can add house elevation in the price equation. Identification comes
from cross-CSO variation in relative elevation within the same sewer
system, not from cross-area elevation differences. 

\subsubsection*{Pumping provisions Soil permeability}

Network distance and elevation together are the two components of
the available hydraulic gradient --- the instrument gains strength
by using both. Additional considerations to strengthen the instrument
include pumping provisions at the time of construction and local sub-catchment
surface imperviousness. A CSO below treatment works elevation requires
pumping, and Victorian designs accounted for this in the construction
of the sewage system. More impervious sub-catchment area upstream
of a CSO results in higher runoff for a given rainfall event and causes
more frequent spilling. However, imperviousness reflects urban density
and land use, which also directly affects house prices (denser areas,
more urban character). To address this, we focus on the geologically-determined
component of imperviousness. In areas with permeable soils (high infiltration
capacity), less rainfall becomes runoff and there is less spilling.
In areas with clay-rich impermeable soils, more rainfall becomes direct
runoff, causing more spilling. Soil permeability is correlated with
Victorian-era development patterns - permeable soils tend to be in
less developed, less dense areas --- and this is absorbed by modern
land use controls. 

We can refine the exposure earlier to add $E_{cw}^{pump}=E_{cw}\times\left(1-Pump_{c}\right)$
and $Pump_{c}\times Tidal_{w}$ that indicates whether Victorian engineering
provided for pumping facilities at the CSO, that is particularly important
for tidal waters, if available.

Further, we can add soil permeability: $E_{cw}^{soil}=E_{cw}\times(1-permeability_{c})/\max_{c'\in w}(1-permeability_{c'})$
were $permeability_{c}$ is the soil class, with clay soils (impermeable)
normalized to receiving value 1. Impermeable soils generate more direct
surface runoff, increasing storm inflow to the combined sewer system.
This is also predetermined because soil type is geological and predates
all human activity. However, soil type correlates with agricultural
land use history, which correlates with when an area was developed,
which correlates with housing stock vintage and quality that need
to be controlled for if using this instrument. The correlation can
be controlled for --- either by including housing age controls or
by interacting soil type with construction era to isolate the geological
component. Soil type enters only through its interaction with rainfall
shocks. The Centre for Ecology and Hydrology HOST classification covers
all areas in Britain.

We are not using the upstream/downstream network position of the CSO.
Network order alone is not a valid instrument because its direction
of effect on spill intensity is theoretically ambiguous. If upstream
CSO weir height is low, upstream CSO fires early and relieves the
interceptor so that the downstream CSO fires later or not at all.
In contrast, if the upstream CSO weir is high the interceptor fills
faster downstream and downstream spill frequency is greater than upstream
spill frequency. However, conditional on the elevation instrument
and WwTW fixed effects, network position could still enter as a control
variable rather than an instrument.

\subsubsection*{WwTW Hydraulics}

The UWWTR dataset gives annual designated Population Equivalent (PE)
--- the standard regulatory measure of sewage load --- for all treatment
works above 2,000 PE in England. Designated PE measures the regulatory
design size of a wastewater treatment works in terms of organic load
(BOD\textsubscript{5}). Because volume of sewage per capita per day
and its BOD concentration are fairly predictable, the designated PE
measures the flow the treatment works was designed to receive in dry
weather --- essentially the engineering benchmark. Designated PE
is largely fixed over time, mostly changing when capital upgrades
are made.\footnote{Alternatively, Flow to Full Treatment (FFT) is a measure of how much
wastewater a treatment works must be able to treat. FFT is normally
calculated as three times the maximum dry weather flow (DWF). Expected
DWF = Consented FFT / 3 and FFT permits are held in the EA's environmental
permit database.} The UWWTR dataset also gives the generated load (actual sewage load
received). It is defined specifically as the maximum average weekly
BOD load entering the treatment plant during the year, excluding unusual
situations such as those due to heavy rain. From the designated and
generated loads, the the fraction of design capacity actually being
used is the loading ratio: 
\[
Stress_{wt}=Actual\,PE_{wt}/Designated\,PE_{wt}
\]
(or its initial mean). When the loading ratio exceeds 1, the works
is overloaded --- CSOs are more likely to spill. When it approaches
1 from below, marginal increases in load have disproportionate effects
on hydraulically constrained CSOs. 

Actual and Designated PE are time-varying and changes to Actual PE
largely reflect rainfall shocks or population/industry growth, while
changes to Designated PE mostly come from upgrades to the system.
In principle, the time variation can be exploited but with additional
controls for population or industrial growth and WwTW-time fixed effects.
Interacting the loading ratio with rainfall shocks then gives a refined
shift measure:
\begin{alignat*}{1}
Stress_{wt}\times\left(Rainfall_{wt}-\mu_{w}\right)/\sigma_{w}
\end{alignat*}
for long-run mean $\mu_{w}$ and standard deviation $\sigma_{w}$.
These storm inflow shocks are common across the catchment (not CSO-specific),
interact differentially with hydraulic components and are plausibly
exogenous to house prices within the catchment, once the rainfall
mean/shock itself is controlled for. 

Rainfall enters the combined sewer system locally through runoff from
impervious surfaces, road drainage, roof connections, and infiltration
into pipes. At each CSO node, inflow therefore depends on rainfall
in its upstream drainage area. However, whether a CSO ultimately activates
depends not only on local inflow but also on system-wide hydraulic
conditions. Once rainfall-induced flows accumulate in the interceptor
network and reach the wastewater treatment works (WwTW), spill activation
depends on whether the works has available treatment capacity. When
the WwTW is near or above its dry-weather design capacity, additional
inflow generates backpressure in the interceptor sewer, increasing
the likelihood that upstream CSOs open to relieve hydraulic load.

This two-stage process implies that spill activation depends on both
local rainfall and catchment-wide hydraulic stress. In practice, because
rainfall is highly spatially correlated within WwTW catchments and
because treatment capacity constraints operate at the works level,
catchment-wide rainfall provides a natural measure of the hydraulic
shock that interacts with predetermined WwTW stress. 

While local rainfall shocks translate into spills more strongly in
stressed systems, catchment rainfall determines the aggregate hydraulic
burden on the treatment works. If spill activation is driven largely
by catchment-wide inflow exceeding WwTW capacity, then the relevant
instrument is catchment rainfall that could be used as an instrument
$R_{wt}\times Stress_{w0}\times Spill_{c0}$. 

\subsubsection*{Unit of Time}

For the unit of time in the shifts, we might use annual or a moving
average of three years. House prices capitalize expected future service
flows, not current conditions. A buyer purchasing near a CSO is discounting
based on their belief about long-run spill frequency --- a belief
formed from cumulative local experience, word of mouth, published
summary statistics, and heuristics --- \textquotedbl this area is
near an old Victorian sewer\textquotedbl{} as a general signal. The
emphasis is on expectations based on experience, not on any single
event. This is directly analogous to the crime literature, where it
is the annual crime rate that is capitalized, not whether a crime
happened on the day of exchange. 

The implication is that annual aggregation is the right frequency,
and a 3-year rolling average is probably better still. Although daily
rainfall and sometimes even daily sewage flows may be available, daily
data adds noise without signal because it introduces enormous day-to-day
variance (a CSO might generate very different flow readings on a wet
November Monday versus a dry July Tuesday) that buyers don't observe
or respond to. The buyer's information set is about cumulated signals
--- the smell from the river last summer, the news story about sewage
spills in the local paper, the EA annual data release.

Annual or three-year averaging also helps with measurement error in
EDM data, since individual spill events may be misrecorded or reflect
equipment failures rather than true hydraulic overflows. But it comes
with one big econometric problem. Annual aggregation reduces identifying
variation and weakens the quasi-experimental rainfall \texttimes{}
stress design by leaving slow-moving variation that may correlate
with neighbourhood trends.

The best approach might be to examine both. Buyers update beliefs
from cumulative recent experience; short-window exposure captures
salience effects, while annual aggregation captures longer-run expectations. 


\end{document}
